\chapter{EUFORIA DISPERATA}\label{euforia-disperata}

\hfill\break

Otto mesi dopo il discorso ufficiale di Wun Ngo Wen all'Assemblea
generale delle Nazioni Unite, le cisterne di colture iperfredde nella
Perielio cominciavano a produrre quantità di carico utile di replicatori
marziani, e a Canaveral e Vandenberg le flotte di razzi Delta VII si
preparavano a lanciarli in orbita. Fu all'incirca in quel periodo che
Wun manifestò il forte desiderio di vedere il Grand Canyon. Ciò che
suscitò il suo interesse fu una vecchia copia di \emph{Arizona Highways}
che qualche patito di biologia aveva lasciato nel suo alloggio.

Me la mostrò un paio di giorni dopo. «Guarda questo» disse, quasi
tremante per l'entusiasmo, mentre ripiegava le pagine di un servizio
fotografico sulla ricostruzione del sentiero di Bright Angel. Il fiume
Colorado che scavava l'arenaria di epoca precambriana formando pozze
verdi. Un turista di Dubai che cavalcava un mulo. «Ne hai mai sentito
parlare, Tyler?»

«Se ho mai sentito parlare del Grand Canyon? Certo. Credo che quasi
tutti lo conoscano.»

«È straordinario. Bellissimo.»

«Spettacolare. Così dicono. Ma non è famoso anche Marte per i suoi
canyon?»

Sorrise. «Ti riferisci alle Fallen Lands. La tua gente le chiamò Valles
Marineris quando le scoprì dall'orbita sessant'anni fa - o centomila
anni fa. Alcune zone assomigliano davvero molto a queste fotografie
dell'Arizona. Ma non ci sono mai stato. E credo che mai ci andrò. Penso,
invece, che mi piacerebbe vedere il Grand Canyon.»

«Allora vallo a vedere. È un paese libero.»

Wun rimase sorpreso dall'espressione - forse era la prima volta che
l'aveva sentita - e fece un cenno con il capo. «Molto bene. Lo farò.
Parlerò con Jason per organizzare il trasporto. Ci verresti anche tu?»

«Dove, in Arizona?»

«Sì! Tyler! In Arizona, al Grand Canyon!» Sarà anche stato un Quarto ma
in quel momento sembrava un bambino di dieci anni. «Ci vieni con me?»

«Devo pensarci.»

Ci stavo ancora pensando quando ricevetti una telefonata da E.D. Lawton.

ooo

Dall'elezione di Preston Lomax, E.D. Lawton era diventato politicamente
invisibile. I suoi contatti nell'industria erano ancora al loro posto -
avrebbe potuto dare una festa ed essere certo che si sarebbe presentata
gente potente - ma non avrebbe mai più esercitato il genere di autorità
ministeriale di cui aveva goduto sotto la presidenza Garland. Anzi,
giravano voci che fosse in uno stato di declino psicologico, rinchiuso
nella sua residenza a Georgetown, a fare telefonate sgradite a ex
alleati politici. Forse era così, ma né Jase né Diane di recente avevano
ricevuto sue notizie; e quando risposi al telefono di casa, rimasi
sbigottito nel sentire la sua voce.

«Vorrei parlarti» esordì.

Il che era interessante, sentendoselo dire dall'uomo che aveva ideato e
finanziato gli atti di spionaggio sessuale di Molly Seagram. Il mio
primo istinto e probabilmente il migliore fu quello di riagganciare. Ma
il gesto mi sembrava inadeguato.

Aggiunse: «Riguarda Jason.»

«Allora parla con Jason.»

«Non posso, Tyler. Non mi ascolterebbe.»

«E questo ti sorprende?»

Sospirò. «Va bene, capisco, sei dalla sua parte, questo è un dato di
fatto. Ma io non voglio ferirlo. Voglio aiutarlo. In effetti è urgente.
È per il suo bene.»

«Non so cosa tu voglia dire.»

«E io non posso dirtelo a questo dannato telefono. Mi trovo in Florida
adesso, sono a venti minuti dall'autostrada. Vieni all'hotel, ti offro
da bere e poi potrai mandarmi a fanculo di persona. Ti prego, Tyler.
Alle 8.00, al bar della hall, sono all'Hilton sulla novantacinquesima.
Forse salverai la vita di qualcuno.»

Riattaccò prima che potessi rispondere.

Chiamai Jason e gli raccontai quello che era accaduto.

«Wow,» disse, e poi «se quello che si dice in giro è vero, passare del
tempo con E.D. dev'essere ancora meno piacevole di quanto lo fosse
prima. Stai attento.»

«Pensavo di non andare all'appuntamento.»

«Non sei certo tenuto a farlo. Ma\ldots{} forse dovresti.»

«Ne ho avuto abbastanza dei giochetti sporchi di E.D., grazie.»

«Però sarebbe meglio sapere cos'ha in mente.»

«Stai dicendo che \emph{vuoi} che lo veda?»

«Solo se te la senti.»

«\emph{Se me la sento}?»

«Sta a te decidere, naturalmente.»

E così salii in macchina e percorsi l'autostrada diligentemente, superai
la bandiera del Giorno dell'Indipendenza (l'indomani sarebbe stato il
quattro) e i venditori di bandiere all'angolo della strada (senza
licenza, pronti a darsela a gambe sui loro pick-up arrugginiti),
ripassando mentalmente tutti i vaffanculo che avevo sempre immaginato di
rivolgere a E.D. Lawton. Quando raggiunsi l'Hilton il Sole era scomparso
dietro i tetti e l'orologio della hall segnava le 8.35.

E.D. era dietro a un separé del bar, e beveva con espressione risoluta.
Sembrò sorpreso di vedermi. Quindi si alzò, mi afferrò il braccio e mi
portò verso la panca in vinile facendomi sedere dall'altra parte del
tavolo di fronte a lui.

«Bevi?»

«Non resterò qui a lungo.»

«Prendi qualcosa, Tyler. Ti farà sentire più bendisposto.»

«Con te ha funzionato? Dimmi solo cosa vuoi, E.D.»

«Capisco se uno è arrabbiato con me dal modo in cui pronuncia il mio
nome, quasi fosse un insulto. Perché sei così incazzato? Per via di
quella cosa con la tua ragazza e il dottore, come si chiama, Malmstein?
Ascolta, ci tengo a dirtelo: non l'ho organizzata io. Non ho neppure
dato il consenso. Quelli che lavoravano per me erano zelanti. Facevano
le cose usando il mio nome. Ecco, volevo che lo sapessi.»

«Questa è una scusa mediocre per un comportamento di merda.»

«Credo di sì. Colpevole. Ti chiedo scusa. Possiamo passare oltre?»

A quel punto sarei potuto uscire. Credo che decisi di restare soltanto
perché emanava un'inquietudine disperata. E.D. era ancora capace di
quella condiscendenza spontanea che lo aveva reso per qualche tempo caro
alla sua famiglia. Ma non era più sicuro di sé. Nei silenzi tra uno
scatto d'ira e l'altro muoveva le mani senza tregua. Si accarezzava il
mento, piegava e riapriva un tovagliolo da cocktail, si lisciava i
capelli. Quel particolare momento di silenzio si protrasse fino a quando
arrivò alla metà del secondo bicchiere. E probabilmente ne aveva già
presi più di due. La cameriera lo aveva servito con disinvolta
familiarità.

«Tu hai una certa influenza su Jason» disse infine.

«Se vuoi parlare con Jason, perché non lo fai direttamente?»

«Perché non posso. Per ovvie ragioni.»

«Allora cosa vuoi che gli dica?»

E.D. mi fissò. Poi guardò il suo bicchiere. «Voglio che tu gli dica di
mettere fine al progetto dei replicatori. Intendo letteralmente. Che
spenga il refrigeratore. Che lo distrugga.»

Adesso toccava a me essere incredulo. «Sai certamente che questo è
improbabile.»

«Non sono stupido, Tyler.»

«Allora perché\ldots»

«È mio figlio.»

«L'hai capito solo ora?»

«Siccome abbiamo avuto discussioni per motivi politici significa che
all'improvviso non è più mio figlio? Credi che sia così superficiale da
non riuscire a fare questa distinzione? Siccome non sono d'accordo con
lui vuol dire che non gli voglio bene?»

«Di te so quello che ho visto.»

«Tu non sai niente.» Cominciò a dire qualcos'altro, ma ci ripensò. «Per
Wun Ngo Wen, Jason è solo una pedina» disse. «Voglio che si svegli e
capisca cosa sta succedendo.»

«Tu l'hai cresciuto perché fosse una pedina. La tua pedina. È solo che
ora non ti piace vedere che qualcun altro ha quella stessa influenza su
di lui.»

«Stronzate. Stronzate. No, va bene, visto che siamo in vena di
confessioni, forse è vero, non lo so, forse per tutti noi ci vorrebbe
una terapia di famiglia, ma non è questo il punto. Il punto è che ogni
persona potente di questo paese finisce per innamorarsi di Wun Ngo Wen e
del suo progetto dei replicatori del cazzo. Per l'ovvia ragione che
costa poco e sembra accettabile per la gente che vota. E chi se ne
importa se non funziona perché \emph{nient'altro} funziona, e se
\emph{niente} funziona allora la fine è vicina e i problemi di tutti
sembreranno diversi quando sorgerà il Sole rosso. Giusto? Non è così? La
infiocchettano, la definiscono una scommessa o un gioco d'azzardo, ma in
realtà è solo un inganno con lo scopo di distrarre i sempliciotti.»

«Analisi interessante,» ribattei, «ma\ldots»

«Starei qui a parlare con te se pensassi che è \emph{un'analisi}
interessante? Fai le domande giuste, se vuoi discutere con me.»

«Del tipo?»

«Tipo, chi è esattamente Wun Ngo Wen? Chi rappresenta e cosa vuole
realmente? Perché, malgrado quello che dicono in televisione, non è
Mahatma Gandhi in un corpo da gnomo. È qui perché vuole qualcosa da noi.
Lo ha voluto fin dal primo giorno.»

«Il lancio dei replicatori.»

«Ovviamente.»

«È un crimine forse?»

«Sarebbe meglio chiedersi: perché i marziani non fanno da sé questo
lancio?»

«Perché non possono prendersi la libertà di parlare a nome dell'intero
sistema solare. Perché un'opera come questa non può essere intrapresa
unilateralmente.»

Alzò gli occhi al cielo. «Sono cose che si \emph{dicono}, Tyler. Parlare
di multilateralismo e diplomazia è come dire ``ti amo''\ldots{} serve
per farsi una scopata. A meno che, certo, i marziani non siano realmente
spiriti angelici discesi dal cielo per liberarci dal male. Cosa che
suppongo tu non creda.»

Wun lo aveva negato così spesso che a stento potei obiettare.

«Voglio dire, guarda la loro tecnologia. Quella gente si è occupata di
biotecnologia ad altissima qualità per mille anni. Se avessero voluto
popolare la galassia con dei nanobot avrebbero potuto farlo molto tempo
fa. Allora perché non l'hanno fatto? Escludendo le spiegazioni che li
vorrebbero migliori di noi, perché? Ovviamente, \emph{perché} temono una
rappresaglia.»

«Una rappresaglia da parte degli Ipotetici? Loro non sanno niente sugli
Ipotetici che noi non sappiamo già.»

«Questo lo dicono loro. Non significa che non li temano. Noi,
invece\ldots{} noi siamo i coglioni che non molto tempo fa hanno
lanciato un attacco nucleare sugli artefatti polari. Sì, ce ne
prenderemo la responsabilità, perché no? Gesù, apri gli occhi, Tyler. È
il classico imbroglio. Non potrebbe essere più falso.»

«O forse sei paranoico.»

«Davvero? A questo punto chi può definire cosa sia la paranoia ai tempi
dello Spin? Siamo tutti paranoici. Noi tutti sappiamo che esistono forze
potenti e maligne che controllano le nostre vite, che è più o meno la
\emph{definizione} di paranoia.»

«Sono solo un medico generico,» dissi, «ma persone intelligenti mi
dicono\ldots»

«Stai parlando di Jason, naturalmente. Jason ti dice che andrà tutto
bene.»

«Non solo Jason. L'intera amministrazione Lomax. Quasi tutto il
Congresso.»

«Ma loro dipendono dal parere dei fanatici. E i fanatici sono
ipnotizzati da tutto questo quanto Jason. Vuoi sapere cosa spinge il tuo
amico Jason? La paura. Ha paura di morire ignorante. Nella situazione in
cui ci troviamo, se lui muore ignorante, significa che la specie umana
muore ignorante. E questo lo spaventa a morte, l'idea che un'intera
specie verosimilmente intelligente possa essere cancellata dall'universo
senza mai capire perché o per cosa. Forse, anziché diagnosticare la mia
paranoia, dovresti pensare alle manie di grandezza di Jason. Comprendere
lo Spin prima di morire è diventata la sua missione. Poi compare questo
Wun e gli dà uno strumento che può usare per quello scopo, e lui
naturalmente lo accetta: è come dare una bustina di fiammiferi a un
piromane.»

«Vuoi davvero che gli dica questo?»

«Io non\ldots» E.D. sembrò improvvisamente cupo, o forse era solo un
picco di alcol nel sangue. «Pensavo, dal momento che lui ti
ascolta\ldots»

«Sai che non è così.»

Chiuse gli occhi. «Forse sì. Non lo so. Ma ci devo provare. Lo capisci?
Per scrupolo di coscienza.» Fui stupefatto che avesse ammesso di averne
una. «Permettimi di essere schietto con te. Mi sento come se stessi
guardando un disastro ferroviario al rallentatore. Le ruote sono uscite
dai binari e il macchinista non se n'è accorto. Allora cosa faccio? È
troppo tardi per tirare l'allarme? Troppo tardi per gridare ``attento''?
Probabilmente sì, Tyler. Ma lui è mio figlio. L'uomo che guida il treno
è mio figlio.»

«Corre gli stessi pericoli che corriamo tutti noi.»

«Credo che ti sbagli. Anche se questa cosa avrà successo, riusciremo a
tirarne fuori solo informazioni astratte. E questo a Jason basta. Ma al
resto del mondo no. Tu non conosci Preston Lomax. Io sì. Lomax sarebbe
più che felice di etichettare Jason come un fallimento e impiccarlo per
questo. Un sacco di gente nella sua amministrazione vuole che la
Perielio venga chiusa o ceduta all'esercito. E queste sono le ipotesi
migliori. Nel caso peggiore, gli Ipotetici perdono la pazienza e
disattivano lo Spin.»

«Ti preoccupa che Lomax chiuda la Perielio?»

«L'ho costruita io la Perielio. Sì, ci tengo. Ma non è questo il motivo
per cui sono qui.»

«Posso riferire a Jason quello che hai detto, ma pensi che cambierà
idea?»

«Io\ldots» A quel punto E.D. si mise a esaminare il piano del tavolo.
Gli occhi si fecero un po' vaghi e lucidi. «No. Ovviamente no. Ma se
vuole parlare\ldots{} voglio che sappia che può mettersi in contatto con
me. Se vuole parlare. Non starò lì a dargli il tormento. Davvero. Solo
se lo vuole.»

Era come se avesse aperto una porta e la sua solitudine interiore avesse
iniziato a fuoriuscire.

Jason credeva che E.D. fosse andato in Florida per prendere parte a un
qualche piano machiavellico. Il vecchio E.D. lo avrebbe fatto. Ma il
nuovo E.D. mi dava l'impressione di un uomo che stava invecchiando,
pieno di rimorsi e di nuovo senza potere, che trovava le sue strategie
nel fondo di un bicchiere e che si era ritrovato in città spinto dal
senso di colpa.

Dissi, in modo più gentile: «Hai provato a parlare con Diane?»

«Diane?» Fece un gesto di sdegno con la mano. «Diane ha cambiato numero.
Non riesco a mettermi in contatto con lei. Comunque, fa parte di quella
dannata setta da fine del mondo.»

«Non è una setta, E.D. È solo una piccola chiesa con delle strane idee.
Simon è più coinvolto di lei.»

«È paralizzata dallo Spin. Proprio come il resto della tua generazione
di merda. Si è gettata a capofitto in questa stronzata religiosa quando
era appena uscita dalla pubertà. Me lo ricordo. Era davvero depressa per
lo Spin. Poi improvvisamente a tavola iniziò a citare Thomas Aquinas.
Volevo che Carol le parlasse. Ma Carol non ne era all'altezza, come al
solito. Allora sai cosa feci? Organizzai un dibattito. Tra lei e Jason.
Da sei mesi non facevano altro che ragionare su Dio. E allora la
trasformai in una cosa ufficiale, come se fosse un torneo di dibattito
studentesco, ed ecco l'artificio, ognuno di loro doveva schierarsi dalla
parte che \emph{non} condivideva: Jason doveva argomentare in favore
dell'esistenza di Dio e Diane doveva assumere il punto di vista
dell'atea.»

Non me ne avevano mai parlato. Ma potevo immaginare con quale sconforto
avessero affrontato i compiti scolastici di E.D.

«Volevo che si rendesse conto di quanto era ingenua. Fece del suo
meglio. Credo che volesse impressionarmi. In pratica ripeteva ciò che
Jason le aveva detto. Ma Jason\ldots» Il suo orgoglio era palese. Gli
brillavano gli occhi e un po' di colore gli ricomparve sul viso. «Jason
fu assolutamente brillante. Brillante e straordinario, meraviglioso.
Jason ribatteva a ogni argomentazione che lei poneva, e molto di più. E
non le ripeteva semplicemente a pappagallo. Aveva studiato teologia,
aveva studiato la cultura biblica. E sorrise per tutto il tempo, come
per dire: ``Guarda, conosco questi argomenti a menadito, li conosco a
fondo quanto te, posso recitarli a occhi chiusi, e continuo a pensare
che siano indegni.'' Fu davvero implacabile. E alla fine del dibattito,
lei pianse. Tenne duro fino alla fine, ma le lacrime le rigavano il
viso.»

Lo fissai.

Notò la mia espressione e fece una smorfia. «Va' all'inferno tu e la tua
superiorità morale. Cercavo di darle una lezione. Volevo che fosse
realista, e non uno di quegli stronzi che guardano solo al loro ombelico
mossi dallo Spin. Tutta la tua generazione del cazzo\ldots»

«T'interessa sapere se è viva?»

«Certo che sì.»

«Nessuno di recente ha avuto sue notizie. Non solo tu, E.D. È
irraggiungibile. Pensavo di provare a rintracciarla. Credi che sia una
buona idea?»

Ma era tornata la cameriera con un altro drink e E.D. stava rapidamente
perdendo interesse per l'argomento, per me, per il mondo attorno a lui.
«Sì, mi piacerebbe sapere se sta bene.» Si tolse gli occhiali e li pulì
con un tovagliolo da cocktail. «Sì, fallo, Tyler.»

E fu così che decisi di accompagnare Wun Ngo Wen nello stato
dell'Arizona.

ooo

Viaggiare con Wun Ngo Wen era come viaggiare con una pop star o un capo
di Stato - molta sicurezza e poca spontaneità, ma parecchio efficiente.
Un susseguirsi strettamente cronometrato di corridoi aeroportuali, voli
charter e scorte in autostrada ci lasciarono, infine, all'inizio del
sentiero di Bright Angel, tre settimane prima del lancio programmato dei
replicatori, in un giorno di luglio, caldo come fuochi d'artificio e
limpido come l'acqua di un ruscello.

Wun se ne stava dove il guardrail seguiva il bordo del canyon. Il
servizio del parco aveva chiuso il sentiero e il centro per i visitatori
ai turisti, e tre dei loro ranger migliori, quelli più fotogenici, erano
pronti a condurre Wun (e un contingente di uomini della sicurezza
federale con fondine a spalla sotto le tenute da escursionismo) in una
spedizione sul fondo del canyon, dove si sarebbero accampati per la
notte.

A Wun era stata promessa riservatezza una volta iniziata l'escursione,
ma ci fu immediatamente una baraonda. I furgoni dei mezzi d'informazione
riempirono l'area del parcheggio, giornalisti e paparazzi si sporgevano
dai cordoni di sicurezza come supplicanti avidi, un elicottero scese in
picchiata facendo riprese lungo il bordo del canyon. Tuttavia Wun era
felice. Sogghignava. Prendeva grandi boccate d'aria dal profumo di pino.
Il caldo era tremendo, soprattutto per un marziano, avrei potuto
pensare, ma lui non mostrava alcun segno di sofferenza malgrado il
sudore che luccicava sulla pelle grinzosa. Indossava una camicia leggera
color cachi in tinta con i pantaloni e un paio di stivali sopra la
caviglia (di quelli da escursionismo, da bambino), che portava ormai da
due settimane per ammorbidirli. Bevve lunghi sorsi da una borraccia
d'alluminio, poi me la offrì.

«Fratello d'acqua» disse.

Risi. «Tienila. Ne avrai bisogno.»

«Tyler, vorrei che facessi la discesa con me. Questo è\ldots» borbottò
qualcosa nella sua lingua, «troppo stufato per un solo pentolone. Troppa
bellezza per un solo essere umano.»

«Puoi sempre condividerlo con gli agenti federali.»

Rivolse un'occhiata malevola ai tizi della sicurezza. «Sfortunatamente
non posso. Loro guardano ma non vedono.»

«Anche questa è un'espressione marziana?»

«Potrebbe anche essere» disse.

ooo

Wun scambiò ancora qualche convenevole con la stampa e il Governatore
dell'Arizona che era appena arrivato, mentre io presi in prestito uno
dei tanti veicoli della Perielio e mi diressi a Phoenix.

Nessuno interferì, nessuno mi seguì; la stampa non era interessata. Sarò
anche stato il medico personale di Wun Ngo Wen - qualcuno degli habitué
della stampa forse mi aveva anche riconosciuto - ma, in assenza di Wun,
non facevo notizia. Neanche lontanamente. Era una bella sensazione.
Aumentai l'aria condizionata tanto che all'interno dell'auto sembrò
essere arrivato l'autunno canadese. Forse si trattava di quella
``euforia disperata'' di cui parlavano i media - il sentimento da
``siamo tutti condannati ma può succedere di tutto'' che aveva raggiunto
l'apice più o meno nel momento in cui Wun si era mostrato pubblicamente.
La fine del mondo e per di più i marziani: considerato questo, cos'era
impossibile? Cos'era anche solo \emph{improbabile}? E da che punto in
poi erano state abbandonate le argomentazioni comuni a favore di decoro,
tolleranza, virtù, cercando di non spezzare l'equilibrio?

E.D. aveva accusato la mia generazione di essere paralizzata dallo Spin,
e forse era vero. Eravamo paralizzati dal terrore ormai da oltre
trent'anni. Nessuno di noi si era mai liberato da quella sensazione di
vulnerabilità sostanziale, quella profonda consapevolezza personale di
avere una spada sospesa sulla testa. Inquinava ogni piacere e faceva
sembrare esitanti e timidi anche i gesti migliori e più coraggiosi.

Ma anche la paralisi si sgretola. Dietro l'inquietudine c'è la
temerarietà. Dietro l'immobilità, l'azione.

Azione non necessariamente buona o saggia, in ogni caso. Passai tre
serie di tabelloni autostradali che mettevano in guardia contro la
possibilità di atti di pirateria ai bordi della strada. Alla radio
locale la cronista per il traffico elencò le strade chiuse per motivi di
pubblica sicurezza, in modo così indifferente da dare l'impressione che
stesse parlando di lavori di manutenzione.

Comunque arrivai senza incidenti all'area di parcheggio dietro il Jordan
Tabernacle.

L'attuale pastore del Jordan era un ragazzo dai capelli a spazzola di
nome Bob Kobel, che aveva accettato per telefono di incontrarmi. Si
avvicinò alla macchina mentre la chiudevo e mi accompagnò nel refettorio
per un caffè, ciambelle e qualche discorso serio. Sembrava un atleta del
liceo un po' ingrassato, ma ancora pieno di quel vecchio spirito di
squadra.

«Ho pensato a quello che mi ha detto» mi disse. «Capisco perché vuole
mettersi in contatto con Diane Lawton. Lei capisce perché è una
questione delicata per questa chiesa?»

«Non proprio, no.»

«La ringrazio per la sua onestà. Lasci che le spieghi, dunque. Sono
diventato pastore di questa congregazione dopo la crisi della giovenca
rossa, ma ne facevo parte già da molti anni. Conosco le persone che la
incuriosiscono\ldots{} Diane e Simon. Un tempo li chiamavo amici.»

«Non più?»

«Mi piacerebbe dire che lo siamo ancora. Ma questo dovrebbe chiederlo a
loro. Vede, dottor Dupree, il Jordan Tabernacle ha avuto una storia
controversa per essere una congregazione relativamente piccola.
Principalmente, perché siamo nati come una chiesa ibrida, un gruppo di
dispensazionalisti di vecchio stampo che si unirono ad alcuni hippy
disillusi del Nuovo Regno. Avevamo in comune una fede risoluta
nell'imminenza della fine dei tempi e un desiderio genuino di fraternità
cristiana. Un'alleanza non facile, come può immaginare. Abbiamo avuto le
nostre controversie. Scismi. Persone che viravano ai margini della
cristianità, dispute dottrinali che, francamente, erano quasi
incomprensibili per molti della congregazione. Ma accadde che Simon e
Diane si schierarono dalla parte di un gruppo di irriducibili
post-tribolazionisti, che rivendicavano come proprio il Jordan
Tabernacle. Questo portò a una gestione difficile delle cose, a quella
che il mondo secolare chiamerebbe lotta di potere.»

«E la persero?»

«Oh, no. Comandarono in modo deciso. Almeno per un po'. Radicalizzarono
il Jordan Tabernacle in un modo che fece sentire moltissimi di noi a
disagio. Dan Condon era uno di loro, e fu lui a coinvolgerci in quella
organizzazione di mine vaganti che cercavano di realizzare la seconda
venuta con una vacca rossa. Cosa che mi sembra ancora grottescamente
presuntuosa. Come se il Signore degli Eserciti potesse aspettare un
programma di allevamento del bestiame prima di chiamare a raccolta i
suoi fedeli.»

Il pastore Kobel sorseggiò il suo caffè.

Dissi: «Non sono in grado di esprimermi sulla loro fede.»

«Lei mi ha detto al telefono che Diane non è più in contatto con la
famiglia.»

«Sì.»

«Potrebbe trattarsi di una sua scelta. Vedevo spesso suo padre alla
televisione. Sembra un uomo che incute timore.»

«Non sono qui per rapirla. Voglio solo assicurarmi che stia bene.»

Un altro sorso di caffè. Un altro sguardo pensieroso.

«Mi piacerebbe dirle che sta bene. E probabilmente è così. Ma dopo gli
scandali, l'intero gruppo si è trasferito in aperta campagna. E alcuni
hanno ancora degli inviti in sospeso da parte di investigatori federali
che vogliono parlare con loro. Quindi le visite sono sconsigliate.»

«Ma non impossibili\ldots»

«Non sono impossibili, se la conoscono. Non sono sicuro che lei possieda
i requisiti, dottor Dupree. Potrei darle delle indicazioni, ma dubito
che la lascerebbero entrare.»

«Anche se lei garantisse per me?»

Il pastore Kobel batté le palpebre. Sembrò pensarci su.

Poi sorrise. Prese un pezzo di carta da un tavolino alle sue spalle e vi
scrisse sopra un indirizzo e qualche riga di indicazioni. «È una buona
idea, dottor Dupree. Lei dica che la manda Padre Bob. Ma stia comunque
attento.»

ooo

Il pastore Kobel mi aveva dato indicazioni per raggiungere la fattoria
di Dan Condon, che risultò essere una casa a due piani in una valle
coperta d'arbusti, a molte ore dalla città. Non un granché come
fattoria, comunque, almeno per i miei occhi inesperti. C'era un grande
fienile, in pessime condizioni rispetto alla casa, e del bestiame che
pascolava tra le chiazze di erba grama.

Appena frenai, un uomo massiccio in salopette scese a balzi i gradini
del portico con i suoi cento chili e passa, la barba folta e
un'espressione infelice. Abbassai il finestrino.

«Proprietà privata, capo» disse.

«Sono qui per vedere Simon e Diane.»

Mi fissò e non disse nulla.

«Non mi stanno aspettando. Ma sanno chi sono.»

«L'hanno invitata? Perché da queste parti non andiamo matti per i
visitatori.»

«Il pastore Bob Kobel ha detto che non vi avrebbe dato fastidio se fossi
passato.»

«Ha detto così, eh?»

«Mi ha detto di riferirvi che sono fondamentalmente innocuo.»

«Il pastore Bob, eh? Ha un documento con sé?»

Presi la carta d'identità, che strinse in mano e portò dentro casa.

Aspettai. Abbassai i finestrini e lasciai che un vento secco soffiasse
all'interno della macchina. Il Sole era talmente basso che le colonne
del portico proiettavano ombre come una meridiana, e quelle ombre si
allungarono un po' di più prima che l'uomo ritornasse, mi restituisse il
documento e dicesse: «Simon e Diane la incontreranno. E mi dispiace se
sono sembrato un po' brusco. Mi chiamo Sorley.» Scesi dalla macchina e
gli strinsi la mano. Aveva una stretta forte. «Aaron Sorley. Ma quasi
tutti mi chiamano fratello Aaron.»

Mi accompagnò all'interno passando per una porta a soffietto rumorosa.
Dentro, la casa era bollente ma vivace. Un bambino con una maglietta di
cotone ci passò accanto di corsa ad altezza ginocchia, ridendo.
Attraversammo una cucina nella quale due donne aiutavano a preparare un
pasto per molte persone - pentoloni sul fuoco, montagne di cavoli sul
tagliere.

«Simon e Diane condividono la camera da letto in fondo, in cima alle
scale, ultima porta a destra\ldots{} può salire.»

Ma non avevo bisogno di una guida. Simon mi stava aspettando in cima
alle scale.

L'ex erede degli scovolini sembrava un po' sciupato. Ma non fu una
sorpresa dato che non l'avevo più visto dalla notte dell'attacco cinese
agli artefatti polari, vent'anni prima. Avrebbe potuto pensare la stessa
cosa di me. Il suo sorriso era ancora straordinario, grande e generoso,
un sorriso che Hollywood avrebbe potuto sfruttare se Simon avesse amato
il denaro più di Dio. Non si accontentò di una stretta di mano. Mi
abbracciò.

«Benvenuto!» disse. «Tyler! Tyler Dupree! Mi scuso se fratello Aaron è
stato un po' brusco poco fa. Non riceviamo molti visitatori, ma
scoprirai che la nostra ospitalità è generosa, perlomeno dopo che hai
varcato la soglia di casa. Ti avremmo invitato prima se avessimo saputo
che c'era una minima possibilità che potessi affrontare il viaggio.»

«Una felice coincidenza» dissi. «Mi trovo in Arizona perché\ldots»

«Oh, lo so. Sentiamo le notizie di tanto in tanto. Sei venuto con l'uomo
grinzoso. Sei il suo medico.»

Mi condusse lungo il corridoio a una porta color crema - la loro porta,
di Simon e Diane - e la aprì.

La stanza al suo interno era arredata in uno stile accogliente, se pur
leggermente antiquato, un grande letto in un angolo con una coperta
imbottita sopra un materasso rigonfio, una tenda di percalle gialla alla
finestra, un tappetino di cotone sul pavimento di tavole chiare. E una
sedia vicino alla finestra. E Diane seduta sulla sedia.

ooo

«È bello vederti» disse. «Grazie per aver trovato del tempo per noi.
Spero che non ti abbiamo distolto dal tuo lavoro.»

«Non più di quanto non volessi esserne distolto. Come state?»

Simon attraversò la stanza e si mise accanto a lei. Poggiò la mano sulla
sua spalla e la tenne lì.

«Stiamo entrambi bene» mi informò lei. «Non ricchi, forse, ma ce la
caviamo. Credo che sia più di quanto chiunque possa aspettarsi di questi
tempi. Mi spiace di non esserci fatti sentire, Tyler. Dopo i problemi al
Jordan Tabernacle è più difficile fidarsi del mondo al di fuori della
chiesa. Immagino che tu abbia sentito parlare di tutto questo\ldots»

«Un casino enorme» intervenne Simon. «La Sicurezza Nazionale ha rimosso
il computer e la fotocopiatrice dalla canonica e li ha portati via senza
più restituirli. Naturalmente, non avevamo niente a che fare con nessuna
di quelle assurdità della giovenca rossa. Ci eravamo limitati a far
circolare qualche opuscolo all'interno della congregazione. Affinché
decidessero, insomma, se era il genere di cose in cui volevano essere
coinvolti. Ecco il motivo per cui siamo stati interrogati dal governo
federale, come puoi immaginare. A quanto pare, questo è un crimine
nell'America di Preston Lomax.»

«Nessun arresto, spero.»

«Nessuno vicino a noi» precisò Simon.

«Ma ha reso tutti nervosi» puntualizzò Diane. «Inizi a pensare a cose
che prima davi per scontate. Telefonate. Lettere.»

«Suppongo che dobbiate essere cauti» dissi.

«Oh, sì» rispose Diane.

«Molto cauti» aggiunse Simon.

Diane indossava una camicia di cotone semplice legata in vita e un
copricapo a quadri rossi e bianchi che sembrava un hijab casereccio.
Niente trucco, ma non ne aveva bisogno. Far indossare a Diane dei
vestiti trasandati era inutile come nascondere un riflettore sotto un
cappello di paglia.

Mi resi conto di quanto avessi avuto voglia anche solo di vederla. Una
voglia assurda. Mi vergognai del piacere che provavo in sua presenza.
Per due decenni eravamo stati poco più che conoscenti. Due persone che
una volta si erano conosciute. Non avevo diritto a quell'aumento del
battito, alla sensazione di accelerazione per l'assenza di gravità che
mi provocava stando semplicemente seduta su quella sedia di legno a
guardarmi e poi girando lo sguardo, arrossendo lievemente quando i
nostri sguardi si incrociavano.

Era irrealistico e ingiusto - ingiusto per qualcuno: forse per me,
probabilmente per lei. Non sarei mai dovuto andare lì.

Lei disse: «E come stai \emph{tu}? Lavori ancora con Jason, immagino.
Spero che stia bene.»

«Sta bene. Manda i suoi saluti.»

Sorrise. «Ne dubito. Non è da Jase.»

«È cambiato.»

«Davvero?»

«Si è parlato molto di Jason,» disse Simon, che continuava a stringerle
la spalla, la mano callosa e scura sul cotone sbiadito. «Di Jason e
dell'uomo grinzoso, il cosiddetto marziano.»

«Non soltanto cosiddetto» lo corressi. «È nato e cresciuto là.»

Simon sbarrò gli occhi. «Se dici così, allora dev'essere vero. Ma come
dicevo, se n'è parlato. La gente sa che l'Anticristo cammina fra noi,
questo è un dato di fatto, e potrebbe pure essere un uomo famoso, in
attesa del momento opportuno, che trama la sua inutile guerra. Pertanto
da queste parti le personalità pubbliche ricevono molti controlli. Non
dico che Wun Ngo Wen sia l'Anticristo, però non sarei l'unico a pensarla
così. Sei a stretto contatto con lui, Tyler?»

«Parlo con lui di tanto in tanto. Non credo sia così ambizioso da poter
essere l'Anticristo.» Anche se E.D. Lawton non sarebbe stato d'accordo
con questa affermazione.

«Tuttavia, questo è il genere di cose che ci porta a essere cauti» disse
Simon. «Ecco perché è stato un problema per Diane rimanere in contatto
con la sua famiglia.»

«Perché Wun Ngo Wen potrebbe essere l'Anticristo?»

«Perché non vogliamo attirare l'attenzione di persone potenti, vista la
prossimità della fine del mondo.»

Non seppi come rispondere a quelle parole.

«Tyler ha viaggiato tanto,» ci interruppe Diane, «forse ha sete.»

Il sorriso di Simon riapparve. «Vorresti qualcosa da bere prima di cena?
Abbiamo un bel po' di gazzosa. Ti piace la Mountain Dew?»

«Sì, va bene».

Uscì dalla stanza. Diane aspettò finché sentimmo i suoi passi per le
scale. Poi alzò la testa e mi guardò in modo più diretto. «Hai fatto
molta strada.»

«Non c'era altro modo per mettermi in contatto.»

«Ma non dovevi darti tanta pena. Sto bene e sono felice. Puoi dirlo a
Jase. E a Carol, per quel che conta. E a E.D., se gli interessa. Non ho
bisogno di una visita di sorveglianza.»

«Non lo è.»

«Ti sei fermato solo per un saluto?»

«A dire il vero, sì, qualcosa del genere.»

«Non ci siamo uniti a una setta. Io non sono sotto minaccia.»

«Non ho detto che tu lo fossi, Diane.»

«Ma lo hai pensato, vero?»

«Sono felice che tu stia bene.»

Girò la testa e la luce del Sole al tramonto attirò il suo sguardo.
«Scusa. Sono solo un po' sorpresa. Vederti così. E sono contenta che ti
vada bene a Est. Te la cavi bene, vero?»

Mi sentii temerario. «No,» dissi, «sono paralizzato. O almeno, questo è
quello che pensa tuo padre. Dice che tutta la nostra generazione è
paralizzata dallo Spin. Siamo tutti ancora bloccati al punto in cui sono
scomparse le stelle. Non lo abbiamo mai accettato.»

«E tu pensi che sia vero?»

«Forse è più vero di quanto ognuno di noi voglia ammettere.» Dicevo cose
che non avevo pensato di tirare fuori. Ma Simon sarebbe tornato a
momenti con la sua lattina di Mountain Dew e il suo sorriso adamantino,
e l'opportunità sarebbe svanita, probabilmente per sempre. «Io ti
guardo,» confessai, «e vedo ancora la ragazza sul prato fuori dalla
Grande Casa. Quindi sì, forse E.D. aveva ragione. Venticinque anni
rubati. Sono trascorsi piuttosto in fretta.»

Diane lo accettò in silenzio. L'aria calda smosse le tende e la stanza
si fece più buia. Poi mormorò: «Chiudi la porta.»

«Non sembrerà strano?»

«Chiudi la porta, Tyler, non voglio che mi sentano.»

Così chiusi piano la porta, lei si alzò, si avvicinò e mi prese le mani.
Le sue erano fresche. «Siamo troppo vicini alla fine del mondo per
mentirci l'un l'altra. Mi dispiace di aver smesso di chiamare, ma ci
sono quattro famiglie che condividono questa casa e un solo telefono, e
risulta piuttosto evidente chi parla con chi.»

«Simon non lo permetterebbe.»

«Al contrario. Simon lo avrebbe accettato. Simon accetta quasi tutte le
mie abitudini e idiosincrasie. Ma non voglio mentirgli. Non voglio
portare questo peso. Ma ammetto che mi sono mancate quelle telefonate,
Tyler. Quelle telefonate erano vita. Quando non avevo denaro, quando la
chiesa si stava dividendo, quando ero triste senza una ragione
valida\ldots{} il suono della tua voce era come una trasfusione.»

«Allora perché smettere?»

«Perché era un atto di slealtà. Allora. E \emph{adesso}.» Scosse la
testa come se stesse cercando di comunicare un concetto difficile ma
importante. «So cosa intendi riguardo allo Spin. Ci penso anch'io. A
volte fingo che esista una realtà in cui non c'è lo Spin e le nostre
vite sono diverse. Le \emph{nostre} vite, la tua e la mia.» Fece un
respiro tremolante e arrossì intensamente. «E se non avessi potuto
vivere in quella realtà, pensavo che avrei potuto almeno visitarla ogni
paio di settimane, chiamarti, essere vecchi amici e parlare di
qualcos'altro oltre alla fine del mondo.»

«E lo consideri sleale?»

«È sleale. Mi sono data a Simon. Simon è mio marito agli occhi di Dio e
della legge. Potrà non essere stata una scelta saggia, ma è stata
comunque la \emph{mia} scelta, e forse non sono il tipo di cristiana che
dovrei essere, ma so cosa significa il dovere, la perseveranza e lo
stare accanto a qualcuno anche se\ldots»

«Anche se \emph{cosa}, Diane?»

«Anche se fa male. Credo che nessuno di noi due abbia bisogno di
esaminare più a fondo le vite che avremmo potuto avere.»

«Non sono venuto qui per renderti triste.»

«No, ma mi fai questo effetto.»

«Allora non mi tratterrò.»

«Ti tratterrai per cena. Solo per educazione.» Mise le mani sui fianchi
e guardò a terra. «Lascia che ti dica una cosa finché abbiamo ancora un
po' di privacy. Per quel che vale. Non condivido tutte le convinzioni di
Simon. Non posso dire onestamente di credere che il mondo finirà con
l'ascesa dei fedeli al cielo. Che Dio mi perdoni, ma non mi sembra
plausibile. Ma credo che il mondo finirà. \emph{Sta} finendo. Sta
finendo da tutta la vita. E\ldots»

La interruppi: «Diane\ldots»

«No, fammi finire. Lasciami confessare. Credo che il mondo finirà. Credo
a quel che Jason mi disse tanti anni fa, che un giorno il Sole si
espanderà e diventerà infernale e in poche ore o giorni il nostro tempo
sulla Terra sarà finito. Non voglio essere sola quel giorno\ldots»

«Nessuno lo vuole.» ``Tranne forse Molly Seagram,'' pensai, ``Molly
Seagram che sente \emph{On the Beach} con il suo flacone di pillole per
suicidarsi. Molly e tutti quelli come lei.''

«E \emph{non} sarò sola. Sarò con Simon. Quello che sto confessando,
Tyler - quello per cui cerco il perdono - è che quando immagino quel
giorno, non è per forza Simon che vedo accanto a me.»

La porta si aprì di colpo. Simon. A mani vuote. «A quanto pare, la cena
è già in tavola,» annunciò, «con una grande brocca di tè freddo per i
viaggiatori assetati. Scendete e unitevi a noi. Ce n'è per tutti.»

«Grazie,» risposi, «sarebbe bello.»

ooo

Gli otto adulti che condividevano la casa erano Dan Condon e sua moglie,
i Sorley, i McIsaac, Simon e Diane. I Sorley avevano tre figli e i
McIsaac cinque, cosicché eravamo in diciassette attorno a un grande
tavolo con i cavalletti nella stanza adiacente alla cucina. Il risultato
fu un piacevole frastuono che durò finché ``zio Dan'' annunciò la
preghiera di ringraziamento, e a quel punto tutte le mani prontamente si
giunsero e tutte le teste prontamente si chinarono.

Dan Condon era il maschio alfa del gruppo. Era alto e quasi sepolcrale,
barba nera, brutto alla Lincoln, e nel benedire il cibo ci ricordò che
dar da mangiare a un forestiero era un atto virtuoso anche se il
forestiero si presentava per caso senza un invito, amen.

Dal modo in cui procedeva la conversazione, conclusi che fratello Aaron
Sorley era comandante in seconda e probabilmente quello che faceva
rispettare le regole quando sorgevano discussioni. Sia Teddy McIsaac che
Simon assecondavano Sorley ma contavano su Condon per i verdetti finali.
La zuppa era troppo salata? «Il giusto» diceva Condon. Il clima caldo
degli ultimi tempi? «Non è insolito in questa parte del mondo»
dichiarava Condon.

Le donne parlavano di rado e per lo più tenevano gli occhi fissi sul
piatto. La moglie di Condon era una donna piccola e tarchiata dai
lineamenti tirati. La moglie di Sorley era grossa quasi quanto lui e
sorrideva visibilmente quando riceveva complimenti per il cibo. La
moglie di McIsaac sembrava avere a malapena diciott'anni in confronto ai
cupi quaranta e più del marito. Nessuna delle donne mi parlò
direttamente, né mi vennero presentate per nome. Diane era un diamante
fra quegli zirconi, saltava all'occhio, e forse questo spiegava il suo
comportamento cauto.

Tutte le famiglie erano rifugiati del Jordan Tabernacle. Non erano i
parrocchiani più radicali, spiegò Zio Dan, come quei dispensazionalisti
dallo sguardo allucinato che erano fuggiti a Saskatchewan l'anno prima,
ma neanche moderati nella loro fede, come il pastore Bob Kobel e la sua
combriccola di gente facile al compromesso. Le famiglie si erano
trasferite alla fattoria (la fattoria di Condon) per allontanarsi di
qualche chilometro dalle tentazioni della città e attendere la chiamata
finale in tranquillità monastica. Fino a quel momento, secondo lui, il
piano aveva avuto successo.

Il resto della conversazione a tavola riguardò un camion con una
batteria scadente, un lavoro di riparazione del tetto ancora in corso e
un'imminente problema della fossa settica. Quando si concluse il pasto
mi sentii sollevato quanto i bambini - Condon rivolse un'occhiata feroce
a una delle figlie dei Sorley quando sospirò troppo forte.

Una volta che i piatti furono ritirati (compito delle donne alla
fattoria di Condon), Simon annunciò che dovevo partire.

Condon mi chiese: «Te la caverai in strada, dottor Dupree? Ormai ci sono
rapine quasi ogni notte.»

«Terrò il finestrino alzato e il pedale dell'acceleratore premuto.»

«Forse sarà meglio.»

Simon disse: «Se non ti dispiace, Tyler, vengo in macchina con te fino
al recinto. Mi piace tornare indietro a piedi, le notti calde come
questa. Anche alla luce di una lanterna.»

Acconsentii.

Poi si misero tutti in fila per un saluto cordiale. I bambini si
dimenarono finché non strinsi loro la mano e li congedai. Quando venne
il suo turno, Diane mi fece un cenno col capo ma abbassò lo sguardo, e
quando le porsi la mano, la strinse senza guardarmi.

ooo

Simon percorse in macchina insieme a me circa quattrocento metri in
salita, agitandosi come chi ha qualcosa da dire ma tiene la bocca
chiusa. Io non lo incalzai. L'aria della notte era fragrante e piuttosto
fresca. Mi fermai dove mi disse lui, in cima a un crinale vicino a una
recinzione rotta e una siepe di ocotillo. «Grazie per il passaggio»
mormorò.

Quando scese si attardò un momento davanti alla portiera aperta.

«Volevi dirmi qualcosa?» chiesi.

Si schiarì la gola. «Sai,» disse alla fine, la voce a malapena più forte
del rumore del vento, «amo Diane tanto quanto amo Dio. Ammetto che
sembra blasfemo. Ho avuto questa sensazione per tanto tempo. Ma credo
che Dio l'abbia mandata sulla Terra per essere mia moglie, che sia il
suo unico scopo. Così ultimamente penso che siano le due facce della
stessa medaglia. Amare lei è il mio modo per amare Dio. Credi che questo
sia possibile, Tyler Dupree?»

Non aspettò una risposta ma chiuse la portiera e accese la torcia
elettrica, e io lo osservai dallo specchietto mentre scendeva con passo
tranquillo la collina in mezzo all'oscurità e al frastuono dei grilli.

ooo

Non m'imbattei in banditi o pirati della strada quella notte.

L'assenza delle stelle e della Luna avevano reso la notte un momento più
oscuro e pericoloso di quanto fosse mai stato, fin dai primi anni dello
Spin. I criminali avevano escogitato delle strategie elaborate per gli
agguati in campagna. Viaggiare di notte aumentava drammaticamente le mie
possibilità di essere rapinato o ucciso.

Per questo motivo il traffico si era diradato durante il viaggio di
ritorno a Phoenix, per lo più camionisti interstatali ben protetti nei
loro autoarticolati a doppio rimorchio. Per gran parte del tempo ero
solo sulla strada, creavo un cuneo luminoso nella notte e ascoltavo lo
stridere delle gomme e l'impeto del vento. Se esiste un suono più
malinconico, non so quale sia. Credo che per questo mettano le radio
nelle macchine.

Ma non c'erano ladri né assassini sulla strada.

Non quella notte.

ooo

E così alloggiai in un motel fuori Flagstaff e il mattino seguente
raggiunsi Wun Ngo Wen e il suo personale di sicurezza nella sala
d'attesa privata dell'aeroporto.

Sul volo per Orlando, Wun era in vena di chiacchiere. Aveva studiato la
geologia del deserto sud-occidentale ed era particolarmente entusiasta
per via di una roccia che aveva acquistato in un chioschetto di souvenir
sulla via del ritorno per Phoenix - costringendo l'intero corteo a
fermarsi e aspettare mentre rovistava in un contenitore di fossili. Mi
mostrò il suo bottino, una concavità calcarea spiraliforme in un grosso
pezzo di argillite di Bright Angel, di circa due centimetri e mezzo di
lato. L'impronta di un trilobite, disse, morto da circa dieci milioni di
anni, recuperata da quei deserti rocciosi e sabbiosi sotto di noi, che
un tempo erano stati il letto di un antico mare.

Non aveva mai visto un fossile prima di allora. Disse che non c'erano
fossili su Marte. Nessun fossile in tutto il sistema solare, tranne che
sull'antica Terra.

ooo

A Orlando ci fecero sedere nei sedili posteriori di un'altra macchina
con un'altra scorta, diretta al complesso della Perielio.

Il Sole stava tramontando, quando ci rimettemmo in strada dopo una
perlustrazione del perimetro che ci aveva tenuto bloccati per un'oretta.
Una volta giunti in autostrada, Wun si scusò per gli sbadigli. «Non sono
abituato a così tanta attività fisica.»

«Ti ho visto sul tapis roulant alla Perielio. Te la cavi bene.»

«Un tapis roulant non è certo un canyon.»

«No, suppongo di no.»

«Ho dolori ovunque ma non sono dispiaciuto. È stata una spedizione
fantastica. Spero che anche tu abbia trascorso il tempo in modo
piacevole.»

Gli dissi che avevo rintracciato Diane e che stava bene.

«Ottimo. Mi dispiace di non averla potuta incontrare. Se è come suo
fratello, deve essere una persona eccezionale.»

«Lo è.»

«Ma la visita non è andata come avevi sperato?»

«Forse speravo nella cosa sbagliata.» Forse avevo sperato nella cosa
sbagliata per molto tempo.

«Be',» disse Wun sbadigliando, con gli occhi socchiusi, «la
domanda\ldots{} come sempre, la domanda è come guardare il Sole senza
esserne accecati.»

Volevo chiedergli cosa volesse dire, ma poggiò la testa contro la
tappezzeria del sedile e mi sembrò più gentile lasciarlo dormire.

ooo

C'erano cinque macchine nel nostro convoglio più un trasporto truppe con
un piccolo reparto di soldati di fanteria in caso di problemi.

Il VTT, veicolo trasporto truppe, era un mezzo squadrato all'incirca
delle dimensioni dei furgoni blindati usati per trasportare denaro da e
verso le banche regionali, e facilmente scambiabile per uno di questi.

Infatti un convoglio della compagnia Brink's si trovò per caso a circa
dieci minuti davanti a noi finché prese l'autostrada verso Palm Bay. Gli
informatori delle bande - collocati sulla strada oltre i principali
snodi e connessi tramite telefono - ci scambiarono per la spedizione
della compagnia Brink's e ci segnalarono come obiettivo di un gruppo di
rapinatori che ci stava aspettando più avanti.

I rapinatori erano criminali sofisticati che avevano già piazzato mine
di superficie in un tratto di strada che costeggiava una riserva
naturale paludosa. Inoltre avevano con sé fucili automatici e un paio di
lanciagranate, e un convoglio della compagnia Brink's non avrebbe avuto
speranze contro di loro: cinque minuti dopo il primo scoppio violento, i
rapinatori si sarebbero dileguati dentro la palude a spartirsi il
bottino. Ma i loro informatori avevano commesso un errore madornale.
Attaccare la spedizione di una banca è una cosa; attaccare cinque
veicoli modificati della sicurezza e un VTT pieno di soldati ben
addestrati e personale di sicurezza è tutt'altra questione.

Tenevo lo sguardo fisso fuori del finestrino oscurato dell'auto,
osservando l'acqua verde in basso e i cipressi spogli che ci passavano
accanto, quando le luci dell'autostrada si spensero.

Un pirata aveva tagliato un cavo elettrico interrato. Improvvisamente
l'oscurità fu davvero intensa, un muro al di là del vetro, niente che
ricambiasse il mio sguardo se non il riflesso della mia espressione
spaventata. Dissi: «Wun\ldots»

Ma lui stava ancora dormendo, la sua faccia rugosa inespressiva come
l'impronta digitale di un pollice.

Poi l'auto di testa investì una mina.

L'urto colpì il nostro veicolo rinforzato come un pugno d'acciaio Il
convoglio rimase prudentemente a distanza, ma noi eravamo abbastanza
vicini da vedere l'auto di testa sollevarsi in una vampata gialla e
ricadere in fiamme sull'asfalto con le ruote spappolate.

Il nostro autista sterzò e, al contrario di ciò che probabilmente gli
era stato insegnato, rallentò. La strada più avanti era bloccata. E poi
ci fu un secondo scoppio dal fondo del convoglio, un'altra mina, che
fece saltare in aria grossi pezzi di asfalto scaraventandoli dentro la
palude e ci bloccò con efficacia implacabile.

Wun era sveglio adesso, confuso e terrorizzato. Aveva gli occhi grandi
come lune e quasi altrettanto luminosi.

Gli spari d'arma da fuoco risuonarono rumorosamente nelle vicinanze. Mi
chinai e tirai Wun verso di me, entrambi piegati in due con le cinture
di sicurezza allacciate di cui forzavamo in maniera convulsa le fibbie.
L'autista si fermò, estrasse un'arma da sotto il cruscotto, e rotolò
fuori dalla portiera.

Nello stesso momento, una decina di uomini dietro di noi si riversò
fuori dal VTT e iniziò a fare fuoco nell'oscurità, cercando di stabilire
un perimetro. Da altri veicoli gli uomini della sicurezza in borghese
cominciarono a convergere verso la nostra auto, cercando di proteggere
Wun, ma gli spari li fermarono prima che riuscissero a raggiungerci.

La rapida risposta doveva avere innervosito i pirati della strada che
quindi aprirono il fuoco con armi pesanti. Uno di loro fece fuoco con
una granata con propulsione a razzo, come mi fu riferito in seguito.
Sapevo solo di essere diventato all'improvviso sordo, l'auto roteava
attorno a uno strano asse e l'aria era piena di fumo e vetri in
frantumi.

ooo

Poi, misteriosamente, mi ritrovai per metà fuori dalla portiera
posteriore, la faccia conficcata sul selciato ghiaioso, il sapore del
sangue in bocca, e Wun vicino a me, poco più in là, disteso sul fianco.
Una delle sue scarpe - gli stivali da bambino per le escursioni che
aveva comprato per la visita al Canyon - aveva preso fuoco.

Lo chiamai per nome. Si mosse, debolmente. I proiettili colpivano i
resti dell'auto dietro di noi, creando fori nell'acciaio. Avevo la gamba
sinistra addormentata. Mi trascinai più vicino a lui e usai un lembo
strappato di tappezzeria per spegnere la scarpa che stava bruciando. Wun
si lamentò e sollevò la testa.

I nostri risposero al fuoco, i proiettili traccianti creavano delle scie
sulle paludi da entrambi i lati della strada.

Wun inarcò la schiena e si sollevò sulle ginocchia. Non sembrava sapere
dove fosse. Perdeva sangue dal naso. Aveva tagli ed escoriazioni sulla
fronte.

«Non alzarti» mi raccomandai con voce roca.

Ma continuò a provare a mettersi in piedi, lo stivale bruciato era a
pezzi ed emanava cattivo odore.

«Per amor di Dio» lo implorai. Allungai la mano verso di lui ma la
scansò. «Per amor di Dio, \emph{non alzarti}!»

Ma alla fine ci riuscì, si sollevò a fatica e si mise in piedi
barcollando, con i rottami in fiamme a fargli da sfondo. Abbassò gli
occhi e sembrò riconoscermi.

«Tyler,» disse, «cos'è successo?»

A quel punto i proiettili lo raggiunsero.

ooo

C'erano molte persone che odiavano Wun Ngo Wen. Diffidavano delle sue
motivazioni, come E.D. Lawton, lo disprezzavano per ragioni più
complesse e meno giustificabili: perché credevano che fosse un nemico di
Dio; perché la sua pelle era casualmente nera; perché abbracciava la
teoria dell'evoluzione; perché incarnava la prova concreta dello Spin e
delle verità inquietanti sull'età dell'universo esteriore.

Molte di quelle persone avevano progettato di ucciderlo. Negli archivi
della Sicurezza Nazionale erano registrate decine di minacce
intercettate.

Ma non venne ucciso da una cospirazione. Lo uccise una combinazione di
avidità, scambio di persona e avventatezza generata dallo Spin.

Fu una morte talmente terrestre da suscitare imbarazzo.

Il corpo di Wun fu cremato (dopo un'autopsia e l'estrazione di campioni
in gran quantità) e venne celebrato in suo onore un funerale di Stato.
Alla cerimonia commemorativa tenuta nella Cattedrale nazionale di
Washington parteciparono dignitari provenienti da tutto il pianeta. Il
Presidente Lomax pronunciò un lungo elogio funebre.

Si parlò di spedire le sue ceneri in orbita, ma non se ne fece mai
niente. Secondo Jason, l'urna fu conservata nel seminterrato dello
Smithsonian Institution in attesa di disposizioni finali.

Probabilmente è ancora là.